\chapter{Introduction}
\label{ch:introduction}

In today's digital world, maintaining a healthy lifestyle has become increasingly challenging. Many users rely on paper notes or memory, which can lead to missed medications or appointments. Health information is usually kept in different places and not saved in one organized system. In many cases, people's memory becomes weaker as they get older. In earlier times, people only arranged their pills physically in boxes. There were no automatic reminders, no data storage, no tracking system, and no proper way to manage overall health information.

Many applications are already available but most of them are not user-friendly, for example: Pill Reminder, Dosecast, Medisafe. These applications provide only medicine reminders; users cannot track their full health routine on one platform. This application brings the trend of mobile-based health management systems along with the problem to patients who missed their medicine dose and other regular routines.

The purpose of this application is to manage health through technology and make a person's medical life easier. As healthcare needs became more complex, users required a system that could integrate multiple types of reminders in one application. Application Smart Healthcare Reminder allows users to maintain medication schedules, record health measurements, manage daily activities, and manage appointments.

A key feature of the system is the integration of doctor and caregiver roles, making it more than just a reminder application. Patients can enter their disease information and send requests to doctors. The doctor can then access the patient's complete health profile. Based on this data, the doctor analyzes the patient's condition and provides personalized recommendations such as medication frequency, measurement schedules, and activity plans.

The doctor can evaluate whether the patient is capable of managing their health independently or requires additional support. If necessary, the system allows the involvement of a caregiver who can manage healthcare activities on behalf of the patient. Depending on the patient's condition:
\begin{itemize}
    \item The patient may manage tasks independently
    \item The caregiver may manage tasks
    \item Or both patient and caregiver can collaboratively manage the healthcare routine
\end{itemize}

The system also includes an AI chatbot that provides intelligent health guidance by responding to user queries, offering basic medical advice, and assisting users in managing their daily healthcare activities. The system will send notifications at scheduled times. The doctor and patient can communicate through their respective dashboards to discuss recommendations and health-related concerns. Emergency contact allows users to save emergency contact details. The patient can quickly access this number in case of an emergency.

The system is designed to be flexible and scalable, which means it can be used by different users with different needs. The application aims to fulfill these needs with interacting features with the help of several tools and technologies along with constraints and limitations.

\section{Problem Statement}
\label{sec:problem-statement}

In today's era, people have a shortage of time, busy routines, or due to multiple health issues, they cannot manage their medical routine. According to studies, between 40\% and 75\% of adults fail to take their prescribed medication on time each day. As a result, individuals fail to take the right dosage at the right time, which may result in additional patient suffering and put their lives in danger.

Many patients cannot manage their daily routine, which is essential for managing their health. There is a lack of coordination between patients, doctors, and caregivers. Doctors do not have access to patient adherence data, making it difficult to analyze patient conditions accurately. Similarly, patients who are critically ill or elderly often require assistance but do not have a proper system to involve caregivers effectively.

This application, Smart Healthcare Reminder, will work as an intelligent mobile-based health assistant which helps patients manage their daily health routine.

\section{Problem Solution}
\label{sec:problem-solution}

Many reminder applications are available in the market but they allow users to enter the medicine name, time, or quantity and give notification at that time. However, this application fulfills this gap and provides all-in-one solutions and not only provides medicine reminders but also manages patients' daily routines such as Measurement (BP, sugar, Weight), Activity (Meal, Water Intake, Walking, Exercise, Yoga, Cycling), Appointment, and Refill.

The Appointment module will help users to remind their medical checkup. The Refill module will help to restock their medicine when stock is limited. The Chatbot will guide users and provide answers to simple health-related questions and allow users to give feedback which can be viewed by admin.

A key enhancement of the system is the integration of doctor and caregiver roles. Patients can send requests to doctors who can then access the patient's complete health profile, including reminder adherence and measurement history. Based on this data, the doctor can analyze the patient's condition and provide personalized recommendations, such as medication adjustments, monitoring frequency, and the need for caregiver support.

After the doctor approves the patient's request, a chat option becomes available for both the doctor and patient, and its caregiver if the doctor gives a recommendation of caregiver to patient. The patient and caregiver can also communicate through their respective dashboards. The doctor and patient can communicate through their respective dashboards to discuss recommendations, health-related concerns, and follow-up information.

\section{Objectives of Proposed System}
\label{sec:objectives}

\begin{enumerate}
    \item The user will authenticate the system with name, email, password, Date of Birth, and gender.
    \item To never miss a dose, allow application ``Medicare'' to send you notifications.
    \item Patient section allows users to send requests, view doctor recommendations, and manage reminders based on their health plan.
    \item Doctor section allows doctors to review patient profiles, analyze health data, and provide recommendations.
    \item After submitting a recommendation, the doctor can provide their Easypaisa or JazzCash number to the relevant patient through the chat. The patient can access the Pay Doctor option under the doctor section, enter the doctor's name and payment number, and submit the transaction ID along with a payment receipt screenshot.
    \item The doctor approves the patient's request, a chat option becomes available for both the doctor and patient.
    \item Caregiver section allows caregivers to monitor patient data, manage health tasks, update reminders, and assist the patient in following their health plan and can also communicate with the patient.
    \item Medication section adds medication to your treatment plan to get reminders.
    \item Measurement section helps to measure health data like blood pressure, weight, and sugar. Setup reminders for daily habits such as walking, drinking water, meal, exercise, yoga, or cycling.
    \item Appointment module will help users to remind their medical checkup.
    \item Pharmacy Medicine Delivery allows patients to order required medicines, check availability, and receive them at their specified delivery address.
    \item Chatbot will guide users and give basic health tips.
    \item Refill module will help to restock their medicine.
    \item Notifications will be sent based on the reminder time set by the user in different modules.
    \item Profile section to view intake progress stored in history section and feedback option will allow users to share their experience and complications.
    \item Emergency contact allows users to save emergency contact details. The patient can quickly access this number in case of an emergency.
\end{enumerate}

\section{Scope}
\label{sec:scope}

The scope of the Smart Healthcare Reminder system is to develop a mobile application that helps users manage their daily health activities in an organized and simple way. Patients can manage their health routines and follow doctor recommendations. Doctors can access patient profiles after appointments, analyze health data, and provide recommendations such as treatment plans, monitoring frequency, and caregiver requirements. Caregivers can assist patients by managing tasks and monitoring their health when needed.

The application will allow users to set reminders for medicines, health measurements (blood pressure, sugar level, weight), physical activities (walking, drinking water, eating meal, exercise, yoga, cycling), medical appointments, and medicine refill. It includes a chatbot that guides users and answers simple health questions. Pharmacy Medicine Delivery allows patients to order required medicines, check availability, and receive them at their specified delivery address. It allows users to give feedback which can be viewed by admin.

It will send notifications at scheduled times and allow users to mark tasks as Done or Dismiss (marked by patient/caregiver). Measurement tracker allows users to record their values (BP (systolic/diastolic), sugar level, weight) and consistently monitor them. The system will also store daily records so users can track their health progress over time. This application is developed for individuals who want to manage their daily health routine and avoid forgetting important health tasks. Emergency contact allows users to save emergency contact details. The main goal of the project is to offer a simple, easy-to-use, and reliable health reminder application where users can manage their health tasks in one place.

\section{System Components}
\label{sec:system-components}

The Application is divided into main parts:
\begin{enumerate}
    \item Patient
    \item Caregiver
    \item Doctor
    \item Pharmacist
    \item Admin
\end{enumerate}

The Application will work in online mode. When the user has an active internet connection, it will allow users to manage Medication, Measurement, Appointments, Activities, Refill, and Profile information. All updates, notifications, and data backup are handled online.

\subsection{Module 1: User Authentication (End User Mobile App)}
The list of features this module will have and how it will be used are the following:
\begin{enumerate}
    \item A splash screen to initialize or load user-specific data if authenticated, else redirected to the authorization screen. The splash screen will check if there's a user logged in or not using Firebase Authentication. If a user is logged in then they will be redirected to the home screen.
    \item The authorization page will have the following options: Log In and Sign Up.
    \item The Sign-Up screen will ask the user to enter personal details like email, full name, date of birth, gender, and password.
    \item The Log In screen will ask the user to enter the registered email and password. After validation, the user will be redirected to the home screen. The login screen will also have an option for forgotten password with email verification.
\end{enumerate}

\subsection{Module 2: Patient}
Patient Module allows users to manage their daily health activities such as setting medication reminders, tracking measurements (BP, sugar, weight), and managing daily routines. Patients can send requests to doctors and follow their recommendations.

\subsection{Module 3: Caregiver}
Caregiver Module allows caregivers to monitor patient health data and manage healthcare tasks on behalf of the patient. Caregivers can view reminders, update task status, and track patient progress.

\subsection{Module 4: Doctor}
Doctor Module allows doctors to view patient profiles after accepting the request and analyze their health data, including measurements and reminder history. Doctors can provide recommendations such as medication plans and monitoring frequency.

\subsection{Module 5: Pharmacist}
The pharmacist receives the request through the Pharmacy Dashboard and provides the medicine that is in stock and delivery cost. The pharmacist then reviews and accepts the delivery address to proceed with the medicine delivery.

\subsection{Module 6: Chat}
The doctor approves the patient's request, a chat option becomes available for both the doctor and patient. The doctor and patient can communicate through their respective dashboards to discuss recommendations and health-related concerns.

\subsection{Module 7: Payment Feature}
After submitting a recommendation, the doctor can provide their Easypaisa or JazzCash number to the relevant patient through the chat. The patient can access the Pay Doctor option under the doctor section, enter the doctor's name and payment number, and submit the transaction ID along with a payment receipt screenshot. The payment remains pending until the doctor reviews it. The doctor can then Approve or Reject the payment.

\subsection{Module 8: Medication}
The user can add a new medication by entering details such as medicine name, dosage, frequency, start date, end date, and reminder time. After saving the information, the medication is added to the reminder list. Once a medication is added, the system automatically gives reminders based on the specified time.

\subsection{Module 9: Measurement}
The user can set reminders for health measurements such as blood pressure, sugar level, and body weight by selecting the measurement type and preferred reminder time.

\subsection{Module 10: Activity}
The user can set reminders for daily activities such as drinking water, eating meals, walking, exercise, yoga, or cycling by selecting an activity type and preferred reminder time.

\subsection{Module 11: Medical Professional}
The Medical Professional module allows patients to connect with doctors and caregivers and manage appointments. Users can send doctor connection requests with their symptoms or invite caregivers using their name and email.

\subsection{Module 12: Refill Medicines}
The user can set refill reminders for their medicines by entering details such as medicine name or stock in the medication section. The system monitors the refill schedule.

\subsection{Module 13: Measurement Tracker}
The user can record and track their health measurements such as blood pressure, blood sugar, and body weight. The system saves the recorded measurements so the user can view them later.

\subsection{Module 14: Pharmacy Delivery Service}
Pharmacy Delivery Service allows patients to search for their required medicines, enter the medicine name and quantity, and check their availability for delivery by the pharmacist.

\subsection{Module 15: Chatbot Assistance}
The chatbot works as an assistant that helps users with basic health-related questions. The user can ask questions about reminders or general health tips.

\subsection{Module 16: Notification and Reminder}
The system automatically generates and schedules notifications based on the reminder time set by the user in different modules. At the scheduled time, the user receives a notification reminding them of a specific task.

\subsection{Module 17: Track Intake}
The track intake module tracks the medicine intake, also measurement (BP, Sugar, weight), daily activities, and appointments.

\subsection{Module 18: Profile Section}
The profile section includes a history option which allows users to view the summary of health-related progress. Users can submit feedback in the feedback option and also include an emergency contact option.

\subsection{Module 19: Admin}
The admin can view registered users. Verify the doctor, pharmacist, and nurse license. Users can submit feedback which can be viewed by the admin.

\section{Related System Analysis / Literature Review}
\label{sec:literature-review}

With the increased use of smartphones and mobile-health technology, many reminder applications have entered the market. Most of the applications currently available primarily focus only on medicine reminders. This proposed application aims to fill this gap by providing an all-in-one solution.

\begin{table}[H]
\centering
\caption{Comparison with Existing Applications}
\label{tab:comparison}
\begin{tabularx}{\textwidth}{|l|X|X|}
\hline
\textbf{Application} & \textbf{Weakness} & \textbf{Proposed Solution} \\ \hline
Medisafe & Provides only medicine reminder. User cannot track their full health routine at one platform. & App manages complete health routine, user will not miss important health tasks. \\ \hline
Dosecast & Gives basic refill reminders, does not give smart options like medicine quantity tracking or low-stock alert. & Medicare tracks medicine and sends a refill reminder when stock is low. \\ \hline
Pill Reminder & Offers a simple alarm-based reminder with no intake tracking or health summary. & Application gives intelligent reminders, intake tracking (Done, Dismiss) and detailed health progress summary. \\ \hline
\end{tabularx}
\end{table}

\section{Vision Statement}
\label{sec:vision}

To create an intelligent mobile application that helps individuals manage their daily health tasks by providing smart reminders for medicine, measurement, activities, appointments, and refill. The application also includes a measurement tracker to record and monitor blood pressure, blood sugar, weight, and a chatbot that provides guidance, answers basic health questions, and helps users manage reminders and coordinate with doctors, caregivers, and pharmacists.

\section{System Limitations and Constraints}
\label{sec:limitations}

\textbf{Limitations:}
\begin{itemize}
    \item \textbf{Internet dependency:} The application needs compatible smartphones and stable internet, which may limit some users.
    \item \textbf{User dependency:} Reminder effectiveness depends on user interaction---correctly setting reminders and responding to notifications.
    \item \textbf{Notification delays:} System or device restrictions can delay notifications.
\end{itemize}

\textbf{Constraints:}
\begin{enumerate}
    \item The application is designed only for Android devices.
    \item The application reminder functionality depends on user allowing notification permission.
\end{enumerate}

\section{Tools and Technologies}
\label{sec:tools}

\begin{table}[H]
\centering
\caption{Tools and Technologies}
\label{tab:tools}
\begin{tabularx}{\textwidth}{|l|l|X|}
\hline
\textbf{Tools/Technology} & \textbf{Version} & \textbf{Rationale} \\ \hline
Android Studio & Panda 2 (2025.2.3) & Official IDE for Android app development. \\ \hline
Firebase Auth & 5.4.1 & For authentication and registration of users. \\ \hline
Draw.io & 24.x & For creating diagrams. \\ \hline
Flutter & 3.38.6 & Familiarity with Flutter; modern apps are developed in Flutter. \\ \hline
GitHub & Built-in Android Studio & Version control and hosting source code. \\ \hline
Cloudinary & Cloud (API) & Cloud storage and file uploads. \\ \hline
Firebase Firestore & 5.6.0 & Database for storing user history, chat, and reminders. \\ \hline
Flutter Local Notification & 17.2.4 & Local notifications to remind or alert the user. \\ \hline
\end{tabularx}
\end{table}

\section{Project Deliverables}
\label{sec:deliverables}

Our system will have the following deliverables by the end of 7th semester:
\begin{enumerate}
    \item Project Proposal
    \item Final Year Project Registration Form
    \item Software Requirement Specification (SRS)
    \item Software Design Specification (SDS)
    \item Prototype
    \item Comprehensive Evaluation Document (Proposal, SRS, SDS)
\end{enumerate}

Our system will have the following deliverables by the end of 8th semester:
\begin{enumerate}
    \item Coding/Implementation
    \item Testing
    \item Deployment
    \item Final
    \item Presentation \& Final Evaluation
\end{enumerate}

\section{Project Planning}
\label{sec:planning}

\begin{table}[H]
\centering
\caption{Project Planning Schedule}
\label{tab:planning}
\begin{tabular}{|c|l|l|l|l|c|}
\hline
\textbf{\#} & \textbf{Task Name} & \textbf{Duration} & \textbf{Start} & \textbf{Finish} & \textbf{Pred.} \\ \hline
1 & Planning & 14 days & 01/12/25 & 14/12/25 & 1 \\ \hline
2 & Requirement Gathering & 17 days & 16/12/25 & 01/01/26 & 2 \\ \hline
3 & Analysis & 30 days & 02/01/26 & 31/01/26 & 2 \\ \hline
4 & Design & 60 days & 02/02/26 & 02/04/26 & 2 \\ \hline
5 & Coding & 125 days & 01/05/26 & 17/09/26 & 2 \\ \hline
6 & Testing & 25 days & 24/08/26 & 18/09/26 & 5 \\ \hline
7 & Implementation & 10 days & 19/09/26 & 28/09/26 & 6 \\ \hline
8 & Documentation & 253 days & 01/12/25 & 30/09/26 & 2 \\ \hline
\end{tabular}
\end{table}

\textbf{Gantt Chart:}\\
\textit{(Include Gantt Chart image here)}\\
% \includegraphics[width=\textwidth]{gantt_chart.png}

\section{Summary}
\label{sec:ch1-summary}

To get together with the emerging trend of mobile-health management systems and solve patients' problems who forget about their dose and regular routine, our project will deliver its modules and fulfill its goals and objectives within the planned timeframe while maintaining the scope of the project to deliver a quality product.
